\documentclass[11pt]{article}

\usepackage[margin=1in]{geometry}
\usepackage{amsmath,amssymb}
\usepackage{booktabs}
\usepackage{enumitem}
\usepackage[T1]{fontenc}
\usepackage[utf8]{inputenc}
\usepackage{xcolor}
\usepackage{hyperref}
\usepackage{microtype}

\hypersetup{
  colorlinks=true,
  linkcolor=blue!55!black,
  citecolor=blue!55!black,
  urlcolor=blue!55!black
}

\newcommand{\software}{\texttt{namaster-proof}}
\newcommand{\Cl}{C_\ell}
\newcommand{\paperVersion}{v2B.0.0}
\newcommand{\paperTimestamp}{2026-07-16 10:25 PT}

\title{\software: Exact pseudo-\(\Cl\) window inference and
tamper-evident provenance for reproducible spin-2 analyses}
\author{Houston Golden}
\date{July 16, 2026 --- \paperVersion}

\begin{document}
\maketitle

\begin{abstract}
\software{} is a focused Python verification layer for two error-prone steps
in cut-sky spin-2 analyses.  First, it evaluates a uniformly rotated
\(EE,EB,BE,BB\) spectrum through the complete NaMaster bandpower-window
operator, avoiding replacement of the operator by bin-centre or
effective-multipole templates.  Second, it publishes JSON results atomically
with content-addressed sidecar receipts and fails closed when either the result
or a declared execution contract changes.  The package also supplies explicit
multipole-support contracts, fixed-grid rotation-angle recovery, command-line
receipt verification, and compatibility tests against the production helpers
from which it was extracted.  The software is intended for method validation
and reproducibility checks; it is not a sky-analysis pipeline, foreground
model, or cosmological inference engine.
\end{abstract}

\section{Summary}

Pseudo-\(\Cl\) estimators make masked spin-2 analyses computationally
practical by relating full-sky spectra to coupled and decoupled
bandpowers~\cite{Hivon2002,Alonso2019}.  Reproducible use of these estimators
requires more than recording a best-fit number: the same window operator must
be used in simulation and inference, harmonic limits must agree across fields
and bins, and intermediate results must remain bound to their execution
metadata.

\software{} packages these checks behind a small public API.  Its numerical
module constructs exact windowed responses for uniform polarization rotation,
evaluates them at one or many angles, recovers an angle on a declared grid,
and compares the window contraction with NaMaster's couple--decouple operator.
Its provenance module atomically publishes a JSON result and a SHA-256-bound
receipt, then verifies the pair and optional metadata.  A third module defines
deterministic multipole and binning contracts.  NumPy is the only required
dependency; a NaMaster workspace is supplied by the user when exercising the
physical window operator.

\section{Statement of Need}

For two spin-2 fields, NaMaster exposes a bandpower-window tensor with spectrum
ordering \([EE,EB,BE,BB]\).  A common shortcut is to evaluate a theory
spectrum at a representative multipole for each bin.  That shortcut is not
generally identical to applying the complete mode-coupling and binning
operator, and a mismatch between the simulation and fitting operators can
appear as a parameter-recovery bias.

A separate reproducibility risk arises in long or sharded calculations.
Ordinary filenames do not prove that a summary corresponds to the intended
seeds, realization count, configuration, or result bytes.  A process can also
terminate between writing a result and its metadata, leaving a plausible but
incomplete artifact.  These problems are usually addressed with analysis-local
helpers.  \software{} provides a reusable, tested implementation with explicit
failure behavior, while remaining narrow enough to audit.

\section{Software Architecture}

The package is divided into three modules:

\begin{description}[leftmargin=2.8cm,style=nextline]
\item[\texttt{windows}]
  Rotation of \(EE,EB,BE,BB\) spectra, exact contraction with a
  \nolinkurl{get_bandpower_windows()} tensor, fixed-grid angle recovery, and
  an equivalence check against \nolinkurl{couple_cell()} followed by
  \nolinkurl{decouple_cell()}.
\item[Multipoles]
  Validation of a common harmonic cutoff and construction of integer
  bandpower edges whose final exclusive edge is \(\ell_{\max}+1\).
\item[\texttt{receipts}]
  Atomic JSON publication, streaming SHA-256 calculation, sidecar-receipt
  construction, and fail-closed content and metadata validation.
\end{description}

The functions accept array-like inputs and return NumPy arrays or plain Python
dictionaries.  NaMaster is deliberately not an installation dependency:
window functions use the small workspace protocol described above, which makes
the algebra testable with synthetic linear operators and permits users to pin
PyMaster independently.

\section{Exact-Window Inference}

For an initially vanishing \(EB\) spectrum and a uniform rotation angle
\(\beta\), the rotated spectra can be decomposed into angle-independent,
\(\cos(4\beta)\), and \(\sin(4\beta)\) components.  In particular,
\begin{align}
 C_\ell^{EE}(\beta)
   &= \tfrac12(C_\ell^{EE}+C_\ell^{BB})
    + \tfrac12(C_\ell^{EE}-C_\ell^{BB})\cos(4\beta),\\
 C_\ell^{BB}(\beta)
   &= \tfrac12(C_\ell^{EE}+C_\ell^{BB})
    - \tfrac12(C_\ell^{EE}-C_\ell^{BB})\cos(4\beta),\\
 C_\ell^{EB}(\beta)=C_\ell^{BE}(\beta)
   &= \tfrac12(C_\ell^{EE}-C_\ell^{BB})\sin(4\beta).
\end{align}
Let \(W_{b\ell}^{ij}\) be the workspace tensor mapping input spectrum \(j\)
to output spectrum \(i\) in band \(b\).  The package precontracts each of the
three components,
\begin{equation}
 R_{b}^{i,k}=\sum_{j,\ell}W_{b\ell}^{ij}C_\ell^{j,k},
\end{equation}
and subsequently evaluates
\begin{equation}
 C_b^i(\beta)=R_b^{i,0}
 +\cos(4\beta)R_b^{i,c}+\sin(4\beta)R_b^{i,s}.
\end{equation}
This retains the full window tensor while making repeated grid evaluation
inexpensive.  The recovery function minimizes a declared weighted or
unweighted squared residual over a user-provided grid; its default is 2001
points spanning \(-1^\circ\) to \(+1^\circ\).

\section{Deterministic Execution and Provenance}

\texttt{publish\_json} serializes a mapping with sorted keys, writes and
flushes a process-specific temporary file, atomically replaces the destination,
and synchronizes the containing directory.  It then publishes a sidecar receipt
containing schema version, result filename, byte count, SHA-256 digest, and
caller-supplied execution metadata.  The content-binding fields are protected
from metadata override.

\texttt{verify\_json\_receipt} requires both files, parses both as JSON
objects, and recomputes every protected field.  \texttt{validate\_json\_receipt}
additionally checks declared metadata such as a suite identifier, seed range,
or realization count.  Missing files, malformed objects, changed bytes,
changed receipts, or contract mismatches raise an exception.  The
\texttt{namaster-proof verify} and \texttt{namaster-proof validate} commands
expose the same behavior to shell workflows and return a nonzero status on
failure.

\section{Validation and Tests}

Version 0.1.0 contains 19 automated tests covering numerical behavior, invalid
inputs, the CLI, receipt mutation, protected fields, and compatibility with
the original production helpers.  The window suite uses an exactly represented
synthetic linear workspace to compare the precontracted response with the
couple--decouple route and to recover an injected grid angle.  Rotation tests
check conservation of total \(EE+BB\) auto-power.  Receipt tests mutate result
bytes and receipt digests independently and require rejection.

Compatibility tests load the pre-package production modules and compare
windowed bandpowers at negative, zero, and positive angles with zero requested
absolute or relative tolerance.  They also compare algebraic spectrum rotation,
bandpower edges, and harmonic keyword contracts.  In the associated physical
production artifact, a workspace of shape
\([4,20,4,1025]\) gave a maximum absolute difference of
\(1.41\times10^{-18}\) between direct window contraction and the
couple--decouple operator.  This number validates that recorded configuration;
it is not a universal error bound.

\section{Worked Examples}

\paragraph{Minimal synthetic operator.}
The documentation constructs an identity-like workspace, builds a response
from synthetic \(EE\) and \(BB\) arrays, evaluates a known rotation, and checks
operator equivalence.  This example has no dependency on a paper-specific
dataset and is suitable for installation tests.

\paragraph{Synthetic CMB recovery campaign.}
The production use case employed CAMB 1.6.6 raw \(EE/BB\) spectra, PyMaster
2.6, \(N_{\rm side}=512\), \(\ell_{\max}=1024\), and deterministic seeds.
For each of two nonzero injected angles, a 500-realization, foreground-free
synthetic run recovered the injection from the mean bandpowers on the declared
\(0.001^\circ\) grid: \(0.270^\circ\mapsto0.270^\circ\) and
\(0.342^\circ\mapsto0.342^\circ\).  The null injection recovered
\(0.000^\circ\).  These are software-recovery checks under the stated
simulation contract, not measurements, detection significances, or evidence
for a physical birefringence model.

\paragraph{Sharded result validation.}
Production shards record fields such as suite, realization count, and seed
range alongside the content digest.  Aggregation can therefore reject a shard
whose bytes were altered or whose declared range does not match the requested
campaign, rather than silently combining it with valid results.

\section{Limitations}

\software{} does not create masks or maps, estimate covariance matrices,
separate cosmic rotation from instrumental polarization calibration, model
foregrounds, choose a likelihood, or infer cosmological parameters.  Uniform
rotation and initially vanishing \(EB\) are assumptions of the optimized
three-component response; the general algebraic rotation helper accepts all
four input spectra, but spatially varying rotation requires another model.
Fixed-grid recovery inherits the range and resolution selected by the caller.
Receipt validation establishes file integrity and agreement with declared
metadata, not the scientific correctness of that metadata.  SHA-256 sidecars
are tamper-evident bindings, not digital signatures or an authorization system.
Schema version 1 has no migration requirement yet; a future incompatible
schema will require an explicit reader or conversion path.

\section{Availability}

Version 0.1.0 is distributed under the MIT License in the
\href{https://github.com/Hubify-Projects/bigbounce}{BigBounce repository}.
Its source is the \software{} package directory.  The package requires Python 3.10 or later
and NumPy 1.24 or later.  Installation, API, CLI, and development commands are
documented in the package README; citation metadata are supplied in
\texttt{CITATION.cff}.  A persistent archival DOI and an independent package
index release are not yet available and must not be inferred from the source
repository.  The complete synthetic-production artifacts used in the worked
example are retained in the same repository's \texttt{reproducibility}
directory, under the \nolinkurl{p1_namaster_500mc} subdirectory.

\section*{Acknowledgements}

The exact-window operations build on NaMaster~\cite{Alonso2019} and the MASTER
pseudo-\(\Cl\) method~\cite{Hivon2002}.  The production example used
CAMB~\cite{Lewis2000} and HEALPix~\cite{Gorski2005}.

\begin{thebibliography}{9}

\bibitem{Hivon2002}
E.~Hivon, K.~M. G\'orski, C.~B. Netterfield, B.~P. Crill, S.~Prunet, and
F.~Hansen, ``MASTER of the CMB anisotropy power spectrum:
A fast method for statistical analysis of large and complex cosmic microwave
background data sets,'' \emph{Astrophysical Journal} 567, 2 (2002),
\href{https://doi.org/10.1086/338126}{doi:10.1086/338126}.

\bibitem{Alonso2019}
D.~Alonso, J.~Sanchez, and A.~Slosar, ``A unified pseudo-\(C_\ell\) framework,''
\emph{Monthly Notices of the Royal Astronomical Society} 484, 4127--4151
(2019), \href{https://doi.org/10.1093/mnras/stz093}{doi:10.1093/mnras/stz093}.

\bibitem{Lewis2000}
A.~Lewis, A.~Challinor, and A.~Lasenby, ``Efficient computation of cosmic
microwave background anisotropies in closed Friedmann--Robertson--Walker
models,'' \emph{Astrophysical Journal} 538, 473--476 (2000),
\href{https://doi.org/10.1086/309179}{doi:10.1086/309179}.

\bibitem{Gorski2005}
K.~M. G\'orski, E.~Hivon, A.~J. Banday, B.~D. Wandelt, F.~K. Hansen,
M.~Reinecke, and M.~Bartelmann, ``HEALPix: A framework for high-resolution
discretization and fast analysis of data distributed on the sphere,''
\emph{Astrophysical Journal} 622, 759--771 (2005),
\href{https://doi.org/10.1086/427976}{doi:10.1086/427976}.

\end{thebibliography}

\end{document}
